\chapter{Sesión 6a: Setup del cliente Flutter y autenticación}
\label{ch:sesion6a}

\section{Objetivo de la sesión}

Crear el cliente móvil multiplataforma en Flutter con:

\begin{itemize}[leftmargin=*]
  \item Estructura de carpetas escalable (core, models, state, screens).
  \item Cliente HTTP con inyección automática del JWT.
  \item Almacenamiento seguro del token en el sistema operativo.
  \item Manejo de estado con \texttt{provider} (\texttt{ChangeNotifier}).
  \item Pantallas de splash, login y registro funcionales contra el backend real.
\end{itemize}

\section{Contexto}

Al cierre de la Sesión~\ref{ch:sesion5}, el backend exponía todos los endpoints necesarios (autenticación, canales, suscripciones, publicaciones, contadores). El siguiente paso natural, según la Fase~6 del protocolo, era construir la aplicación móvil multiplataforma con Flutter que consumiera esa API.

Se optó por dividir la Sesión~6 en cuatro sub-sesiones para que cada una entregue algo verificable:

\begin{enumerate}[leftmargin=*]
  \item \textbf{6a (este capítulo):} setup + autenticación.
  \item \textbf{6b:} listado, búsqueda, detalle y creación de canales.
  \item \textbf{6c:} publicaciones dentro del canal y contadores de no leídos.
  \item \textbf{6d:} notificaciones push reales con Firebase Cloud Messaging.
\end{enumerate}

\section{Desarrollo}

\subsection{Generación del proyecto}

Se creó el proyecto Flutter con organización \texttt{mx.uam}:

\begin{lstlisting}[language=bash]
cd ~/Escritorio/notiuam
flutter create mobile --org mx.uam --project-name notiuam_mobile
\end{lstlisting}

\subsection{Dependencias}

Se configuraron en \texttt{pubspec.yaml}:

\begin{lstlisting}[language=yaml]
dependencies:
  flutter:
    sdk: flutter
  cupertino_icons: ^1.0.8
  dio: ^5.7.0                       # Cliente HTTP
  flutter_secure_storage: ^9.2.2    # Token seguro en Keychain / EncryptedSharedPreferences
  provider: ^6.1.2                  # Manejo de estado
  intl: ^0.19.0                     # Formato de fechas

dev_dependencies:
  flutter_test:
    sdk: flutter
  flutter_lints: ^5.0.0
\end{lstlisting}

\subsection{Estructura de carpetas}

\begin{lstlisting}[language=bash]
mobile/lib/
|- main.dart                          (punto de entrada, MultiProvider)
|- core/
|  |- constants.dart                  (ApiConfig: baseUrl segun plataforma)
|  |- token_storage.dart              (guardado seguro del JWT)
|  `- api_client.dart                 (Dio + interceptor de autenticacion)
|- models/
|  |- user.dart
|  |- auth_response.dart
|  |- channel.dart
|  `- post.dart
|- state/
|  `- auth_provider.dart              (ChangeNotifier de autenticacion)
`- screens/
   |- splash_screen.dart
   |- login_screen.dart
   |- register_screen.dart
   `- home_shell.dart                 (barra de navegacion inferior)
\end{lstlisting}

\subsection{Decisiones técnicas}

\begin{itemize}[leftmargin=*]
  \item \textbf{\texttt{dio} en lugar de \texttt{http}.} Dio ofrece interceptores, manejo estructurado de errores (\texttt{DioException}) y mejor configuración de timeouts.
  \item \textbf{Interceptor de autenticación.} El \texttt{ApiClient} lee el token desde \texttt{TokenStorage} en cada petición y lo añade al header \texttt{Authorization: Bearer <token>} automáticamente. Las pantallas no se preocupan por el token.
  \item \textbf{\texttt{flutter\_secure\_storage} en lugar de \texttt{shared\_preferences}.} Guarda el JWT en el Keychain de iOS o en \texttt{EncryptedSharedPreferences} de Android. Nunca en texto plano.
  \item \textbf{\texttt{provider} con \texttt{ChangeNotifier}.} Patrón estándar de la comunidad Flutter. Alternativa a \texttt{bloc} que se descartó por ser más verboso para el tamaño del proyecto.
  \item \textbf{\texttt{ProxyProvider} para dependencias encadenadas.} El \texttt{ApiClient} depende del \texttt{TokenStorage}, y el \texttt{AuthProvider} del \texttt{ApiClient}. Se usan \texttt{ProxyProvider} y \texttt{ChangeNotifierProxyProvider} para inyectarlos sin duplicar instancias.
\end{itemize}

\subsection{Detección de plataforma para la URL del backend}

El emulador de Android usa \texttt{10.0.2.2} para referirse al host, mientras que Linux/iOS usan \texttt{localhost}. El archivo \texttt{constants.dart} encapsula la lógica:

\begin{lstlisting}[language=Java]
class ApiConfig {
  static String get baseUrl {
    if (Platform.isAndroid) {
      return 'http://10.0.2.2:8081';
    }
    return 'http://localhost:8081';
  }
  static const Duration timeout = Duration(seconds: 15);
}
\end{lstlisting}

\subsection{Configuración Android}

Se editó \texttt{android/app/src/main/AndroidManifest.xml} para:

\begin{itemize}[leftmargin=*]
  \item Añadir el permiso \texttt{android.permission.INTERNET}.
  \item Habilitar \texttt{android:usesCleartextTraffic="true"} (necesario porque el backend sirve HTTP en desarrollo, no HTTPS).
\end{itemize}

\begin{aviso}[Advertencia]
\texttt{usesCleartextTraffic="true"} solo es aceptable en \texttt{debug}. En la versión de producción habrá que exponer el backend por HTTPS con un certificado válido.
\end{aviso}

\subsection{Cambios a archivos previos}

Ninguno. Esta es una sub-sesión completamente nueva.

\subsection{Problemas encontrados y solución}

\subsubsection{Falta el SDK de Android}

Al primer \texttt{flutter run} apareció:

\begin{lstlisting}[language=bash]
SDK location not found. Define a valid SDK location with an
ANDROID_HOME environment variable or by setting the sdk.dir path...
\end{lstlisting}

\textbf{Solución:} instalar Android Studio por Snap (\texttt{sudo snap install android-studio --classic}), configurar el SDK siguiendo el asistente, y añadir las variables al shell:

\begin{lstlisting}[language=bash]
export ANDROID_HOME=$HOME/Android/Sdk
export PATH=$PATH:$ANDROID_HOME/platform-tools
export PATH=$PATH:$ANDROID_HOME/cmdline-tools/latest/bin
\end{lstlisting}

\subsubsection{Clang++ no encontraba \texttt{libstdc++}}

Al intentar compilar para Linux desktop:

\begin{lstlisting}[language=bash]
/usr/bin/ld: no se puede encontrar -lstdc++: No existe el archivo o el directorio
\end{lstlisting}

El diagnóstico con \texttt{clang++ -v} reveló que clang usa por defecto GCC~14 (el más reciente) mientras que \texttt{g++} apunta a GCC~13. Faltaba el paquete \texttt{libstdc++-14-dev}.

\textbf{Solución:}

\begin{lstlisting}[language=bash]
sudo apt install -y libstdc++-14-dev
\end{lstlisting}

\subsubsection{Falta \texttt{libsecret-1-dev}}

El plugin \texttt{flutter\_secure\_storage\_linux} requiere \texttt{libsecret} para acceder al llavero del sistema:

\begin{lstlisting}[language=bash]
The following required packages were not found:
  - libsecret-1>=0.18.4
\end{lstlisting}

\textbf{Solución:}

\begin{lstlisting}[language=bash]
sudo apt install -y libsecret-1-dev
\end{lstlisting}

\subsubsection{Archivos \texttt{api\_client.dart} y \texttt{token\_storage.dart} con contenido duplicado}

Ambos archivos contenían el mismo bloque de \texttt{constants.dart} por un error de copiado. La compilación falló con errores de tipo \texttt{'ApiClient' isn't a type} y \texttt{'TokenStorage' isn't a type}.

\textbf{Solución:} sobrescribir ambos archivos con el contenido correcto mediante una escritura completa del archivo para garantizar que no quedaran partes duplicadas.

\section{Verificación}

\subsection{Arranque}

\begin{lstlisting}[language=bash]
cd ~/Escritorio/notiuam/mobile
flutter clean
flutter run -d linux
\end{lstlisting}

La primera compilación tarda 1--2 minutos. Al terminar se abre una ventana nativa con la app NotiUAM.

\subsection{Flujo verificado}

\begin{enumerate}[leftmargin=*]
  \item \textbf{Splash} visible 1 segundo mientras \texttt{AuthProvider.bootstrap()} verifica si hay token guardado.
  \item \textbf{Login} con \texttt{aaron@uam.mx} / \texttt{password123} → entra al \texttt{HomeShell}.
  \item \textbf{HomeShell} con 4 pestañas: Explorar, Mis canales, Suscritos, Perfil.
  \item \textbf{Logout} desde el AppBar → vuelve al login y limpia el token del storage.
  \item \textbf{Persistencia:} cerrar la app y reabrirla mantiene la sesión activa (token en Keychain / EncryptedSharedPreferences).
\end{enumerate}

\section{Resultado}

\begin{exito}[Estado al cierre]
\begin{itemize}[leftmargin=*]
  \item Proyecto Flutter funcional en Linux desktop (y listo para Android).
  \item Cliente HTTP con JWT automático.
  \item Almacenamiento seguro del token.
  \item Login y registro funcionales contra el backend real.
  \item Arquitectura de estado escalable para las próximas sub-sesiones.
\end{itemize}
\end{exito}

\textbf{Pendiente para la siguiente sub-sesión:} reemplazar los placeholders del \texttt{HomeShell} por pantallas reales de canales.